\subsection{A quantitative extract-and-restore lemma}\label{sec:global:fork}
Put $a_\kappa:=2^{-\kappa}$.
The central improvement over Chapter~11 is to
\emph{condition on the already sampled accepting
main path} and search only for an \emph{alternative}
accepted inner challenge string. The main path is
never replaced by the successful fork.
The elementary probability bound below makes
this strategy work uniformly over all adaptive
pre-block histories.

\begin{definition}[Local accepted challenge tree]\label{def:global:tree}
Fix a pre-block public history, a public OR
statement $(C,S)$, and the entire saved private
state and random tape of a causal prover.
For each challenge string
$e\in\{0,1\}^{\kappa}$,
run the deterministic prover from that state.
Let
\[
 \mathcal A:=\{e:\text{all }\kappa
 \text{ mini-proofs accept}\},\quad
 m:=|\mathcal A|,\quad
 p:=m/2^\kappa.
\]
For an accepted main string $e\in\mathcal A$,
a side run is a \emph{fork} if it accepts on
some string $e'\in\mathcal A\setminus\{e\}$.
\end{definition}

\begin{lemma}[Probability of a distinct accepting side string]\label{lem:global:side-success}
For every fixed accepted main string $e$,
one independently sampled uniform side
challenge string gives a fork with probability
\[
 \boxed{p-a_\kappa}.
\]
\end{lemma}
\begin{proof}
The accepted challenge set has $m$ elements,
exactly one of which is the fixed main
string $e$. Each of the other $m-1$
strings has probability $2^{-\kappa}$.
Thus the success probability equals
$(m-1)2^{-\kappa}
=p-2^{-\kappa}=p-a_\kappa$.
\end{proof}

\begin{lemma}[Exact failure probability after bounded replay]\label{lem:global:exact-fail}
Fix the deterministic pre-block state.
Run a fresh uniform main OR block and,
only if it accepts, make $R\ge1$ independent
side attempts for a \emph{distinct}
accepted challenge string.
The joint probability that the main block
accepts but no fork is found is
\[
 \boxed{p(1-p+a_\kappa)^R.}
\]
\end{lemma}
\begin{proof}
The main string is accepted with
probability $p$.
Conditioned on any accepted main
string, Lemma~\ref{lem:global:side-success}
makes the failure probability of one
side attempt $1-p+a_\kappa$.
The side challenge strings are mutually
independent and also independent of the
main sample. Thus all $R$ attempts
fail with probability
$(1-p+a_\kappa)^R$.
This expression does not depend on
the particular accepted main string.
Multiplying by $p$ gives the formula.
\end{proof}

\begin{lemma}[Uniform elementary replay bound]\label{lem:global:uniform-fail}
For every admissible
$p=m2^{-\kappa}$ and every integer
$R\ge1$,
\[
 \boxed{
 p(1-p+a_\kappa)^R
 \le a_\kappa+\frac{1}{R+1}.}
\]
\end{lemma}
\begin{proof}
If $m=0$, the left side is zero.
Otherwise $p\ge a_\kappa$.
Put $u:=p-a_\kappa\in[0,1]$.
Then
\[
 p(1-p+a_\kappa)^R
 =(u+a_\kappa)(1-u)^R
 \le u(1-u)^R+a_\kappa.
\]
The real-valued function
$f(u)=u(1-u)^R$ on $[0,1]$
is zero at the endpoints.
For $0<u<1$ its derivative is
\[
 f'(u)=(1-u)^{R-1}(1-(R+1)u).
\]
Hence the maximum is attained at
$u=1/(R+1)$ and equals
$R^R/(R+1)^{R+1}$,
which is at most $1/(R+1)$.
Substitution gives the required bound.
\end{proof}

\begin{theorem}[Per-block online extraction with a preserved main path]\label{thm:global:one-block}
Fix a public OR statement before its
first one-bit challenge and let a causal
randomized classical prover be
checkpoint-replayable.
There is a black-box procedure that:
\begin{enumerate}[label=\textup{(\roman*)}]
\item runs the unmodified main OR block
against $\kappa$ independent uniform
challenge bits;
\item if it accepts, makes at most $R$
independent side replays and attempts
the explicit special-soundness
extraction of Chapter~11;
\item restores the exact main post-block
state before the next protocol stage.
\end{enumerate}
It produces a valid OR opening
on every successful fork, and
\[
 \boxed{
 \Pr[\text{main accepts but extraction fails}]
 \le 2^{-\kappa}+\frac{1}{R+1}.}
\]
The bound holds for every pre-block
history, including adaptively
chosen public statements.
\end{theorem}
\begin{proof}
Condition on the complete saved
prover state and tape at the start
of the block. The statement is
then fixed and
Definition~\ref{def:global:tree}
gives the deterministic accepted
challenge set.
By Lemma~\ref{lem:global:uniform-fail},
the probability that main acceptance
is followed by $R$ failed side
attempts is at most
$2^{-\kappa}+1/(R+1)$.

If a distinct accepting side string
is found, the two runs start from
the identical saved state and tape.
Lemma~\ref{lem:mv:earliest}
supplies two accepting transcripts
with the same OR first message
and different one-bit challenges
at the earliest differing index.
Theorem~\ref{thm:mv:two-run-extract}
computes a valid permitted opening.
The procedure then restores the main
post-block state by
Definition~\ref{def:global:access}.
The conditional bound is uniform
over saved states, so averaging
over the prover tape and all
earlier adaptive public histories
preserves it.
\end{proof}

\begin{corollary}[Uniform extraction error over sequential blocks]\label{cor:global:J-blocks}
For at most $J$ sequential OR blocks,
suppose every accepted block is followed
by the extract-and-restore procedure of
Theorem~\ref{thm:global:one-block}.
The probability that the \emph{main
protocol accepts} yet at least one
accepted OR block fails to yield an
extracted opening is at most
\[
 \boxed{
 \varepsilon_{\rm fork}(J,\kappa,R)
 :=\min\left\{1,
 J\left(2^{-\kappa}+\frac1{R+1}\right)\right\}.}
\]
\end{corollary}
\begin{proof}
At the start of the $j$-th reached
block, condition on the entire
history and saved prover state.
Theorem~\ref{thm:global:one-block}
bounds the event that this block
accepts but extraction fails.
If the final main protocol accepts,
every reached mandatory OR block
accepted, so failure at any such
block lies in one of these events.
A union bound over at most $J$
block positions gives the displayed
bound, capped by one.
\end{proof}

\begin{lemma}[Expected and worst-case replay overhead]\label{lem:global:runtime}
The $J$-block procedure uses at most
$JR$ side block executions,
and at most $J(R+1)$ total
main-plus-side block executions.
If a block contains $\kappa$
one-bit mini-proofs, the number
of executed mini-proof instances
is at most $\kappa J(R+1)$.
\end{lemma}
\begin{proof}
At each reached block the procedure
makes exactly one main execution
and at most $R$ side executions.
There are at most $J$ blocks.
Multiply the resulting number
of blocks by the maximum of
$\kappa$ mini-proofs per block.
\end{proof}

\begin{remark}[Why this is stronger than isolated forking]
Chapter~11 obtained two accepting
runs for one fixed statement.
Theorem~\ref{thm:global:one-block}
instead \emph{retains the actual main
accepted run}, constructs a fork only
as a private side computation, and
restores the precise state that will
receive the next outer challenge.
The failure bound is unconditional
over histories; no lower bound on
the prover's conditional acceptance
probability needs to be assumed.
\end{remark}
